% Lokale semantische Schlussformen. Keine neuen Mengenaxiome.
\subsection*{Einzeln referenzierte semantische Schlussformen}

In diesem Abschnitt wird die durch
\FormulaRefAuto{B48SatisfactionExists}[thm] konstruierte und eindeutig
bestimmte Erfuellungsrelation verwendet. Die definierenden Klauseln
aus \FormulaRefAuto{B48Satisfaction}[def] sind durch die dortigen
Proofparts gerechtfertigt. Die folgenden Saetze zeigen, welche
lokalen Schluesse aus diesen Klauseln folgen.
Alle Zeilen sind Aussagen der untersuchenden Mengenlehre.
Insbesondere sind \(M,s\models\neg\varphi\) und
\(\neg(M,s\models\varphi)\) verschiedene Schreibweisen, deren
Zusammenhang jeweils durch die angegebene Definition begruendet wird.

\FormulaThmDeltaK[Keine Belegung erfuellt die Widerspruchsformel]
{\neg(M,s\models\bot)}
{B48SatNotBottom}{
 \DeltaRow{Struktur und Belegung}{M=(D,E),\quad s:\operatorname{Var}\to D}
}
\begin{tabproof}
 \proofstep{1}{M,s\models\bot}{\rA}
 \proofstep{1}{\bot}{\FormulaRefAuto{B48Satisfaction}{1}}
 \proofstep{}{\neg(M,s\models\bot)}{\rCI{1,2}}
\end{tabproof}

\FormulaThmDeltaK[Semantische Konjunktionseinfuehrung]
{M,s\models\varphi,\ M,s\models\psi
 \ \vdash\ M,s\models(\varphi\land\psi)}
{B48SatAndIntro}{
 \DeltaRow{Formeln}{\varphi,\psi}
 \DeltaRow{Belegung}{s:\operatorname{Var}\to D}
}
\begin{tabproof}
 \proofstep{1}{M,s\models\varphi}{\rA}
 \proofstep{2}{M,s\models\psi}{\rA}
 \proofstep{1,2}{(M,s\models\varphi)\land(M,s\models\psi)}
 {\rAI{1,2}}
 \proofstep{1,2}{M,s\models(\varphi\land\psi)}
 {\FormulaRefAuto{B48Satisfaction}{3}}
\end{tabproof}

\FormulaThmDeltaK[Erste semantische Konjunktionsbeseitigung]
{M,s\models(\varphi\land\psi)\ \vdash\ M,s\models\varphi}
{B48SatAndLeft}{
 \DeltaRow{Formeln}{\varphi,\psi}
 \DeltaRow{Belegung}{s:\operatorname{Var}\to D}
}
\begin{tabproof}
 \proofstep{1}{M,s\models(\varphi\land\psi)}{\rA}
 \proofstep{1}{(M,s\models\varphi)\land(M,s\models\psi)}
 {\FormulaRefAuto{B48Satisfaction}{1}}
 \proofstep{1}{M,s\models\varphi}{\rAEa{2}}
\end{tabproof}

\FormulaThmDeltaK[Zweite semantische Konjunktionsbeseitigung]
{M,s\models(\varphi\land\psi)\ \vdash\ M,s\models\psi}
{B48SatAndRight}{
 \DeltaRow{Formeln}{\varphi,\psi}
 \DeltaRow{Belegung}{s:\operatorname{Var}\to D}
}
\begin{tabproof}
 \proofstep{1}{M,s\models(\varphi\land\psi)}{\rA}
 \proofstep{1}{(M,s\models\varphi)\land(M,s\models\psi)}
 {\FormulaRefAuto{B48Satisfaction}{1}}
 \proofstep{1}{M,s\models\psi}{\rAEb{2}}
\end{tabproof}

\FormulaThmDeltaK[Unvertraeglichkeit einer Formel mit ihrer Negation]
{M,s\models\varphi,\ M,s\models\neg\varphi\ \vdash\ \bot}
{B48SatContradiction}{
 \DeltaRow{Formel}{\varphi}
 \DeltaRow{Belegung}{s:\operatorname{Var}\to D}
}
\begin{tabproof}
 \proofstep{1}{M,s\models\varphi}{\rA}
 \proofstep{2}{M,s\models\neg\varphi}{\rA}
 \proofstep{2}{\neg(M,s\models\varphi)}
 {\FormulaRefAuto{B48Satisfaction}{2}}
 \proofstep{1,2}{\bot}{\rBI{1,3}}
\end{tabproof}

\FormulaThmDeltaK[Semantischer Negationsschluss]
{\bigl((M,s\models\varphi)\to\bot\bigr)
 \ \vdash\ M,s\models\neg\varphi}
{B48SatNegationIntro}{
 \DeltaRow{Formel}{\varphi}
 \DeltaRow{Belegung}{s:\operatorname{Var}\to D}
}
\begin{tabproof}
 \proofstep{1}{(M,s\models\varphi)\to\bot}{\rA}
 \proofstep{2}{M,s\models\varphi}{\rA}
 \proofstep{1,2}{\bot}{\rRE{1,2}}
 \proofstep{1}{\neg(M,s\models\varphi)}{\rCI{2,3}}
 \proofstep{1}{M,s\models\neg\varphi}
 {\FormulaRefAuto{B48Satisfaction}{4}}
\end{tabproof}

\FormulaThmDeltaK[Semantischer klassischer Widerspruchsschluss]
{\bigl((M,s\models\neg\varphi)\to\bot\bigr)
 \ \vdash\ M,s\models\varphi}
{B48SatClassical}{
 \DeltaRow{Formel}{\varphi}
 \DeltaRow{Belegung}{s:\operatorname{Var}\to D}
}
\begin{tabproof}
 \proofstep{1}{(M,s\models\neg\varphi)\to\bot}{\rA}
 \proofstep{2}{\neg(M,s\models\varphi)}{\rA}
 \proofstep{2}{M,s\models\neg\varphi}
 {\FormulaRefAuto{B48Satisfaction}{2}}
 \proofstep{1,2}{\bot}{\rRE{1,3}}
 \proofstep{1}{M,s\models\varphi}{\rCE{2,4}}
\end{tabproof}

\FormulaThmDeltaK[Semantische Existenzeinfuehrung mit bestimmtem Zeugen]
{a\in D,\ M,s[x\mapsto a]\models\varphi
 \ \vdash\ M,s\models\exists x\,\varphi}
{B48SatExistsIntro}{
 \DeltaRow{Formel und Variable}{\varphi,\quad x\in\operatorname{Var}}
 \DeltaRow{Belegung}{s:\operatorname{Var}\to D}
}
\begin{tabproof}
 \proofstep{1}{a\in D}{\rA}
 \proofstep{2}{M,s[x\mapsto a]\models\varphi}{\rA}
 \proofstep{1,2}{a\in D\land(M,s[x\mapsto a]\models\varphi)}
 {\rAI{1,2}}
 \proofstep{1,2}{\exists b\,
   (b\in D\land(M,s[x\mapsto b]\models\varphi))}
 {\rEI{3}}
 \proofstep{1,2}{M,s\models\exists x\,\varphi}
 {\FormulaRefAuto{B48Satisfaction}{4}}
\end{tabproof}

\FormulaThmDeltaK[Semantische Existenzbeseitigung mit Zeugenentladung]
{\begin{gathered}
 M,s\models\exists x\,\varphi,\\
 \forall a\,(a\in D\to
       ((M,s[x\mapsto a]\models\varphi)\to Q))
 \ \vdash\ Q
 \end{gathered}}
{B48SatExistsElim}{
 \DeltaRow{Formel und Variable}{\varphi,\quad x\in\operatorname{Var}}
 \DeltaRow{Belegung}{s:\operatorname{Var}\to D}
 \DeltaRow{Aussage der untersuchenden Theorie}{Q}
 \DeltaRow{Frischer Zeuge}{b}
}
Der Zeuge \(b\) kommt in \(Q,M,D,s,\varphi\) und in den ausserhalb
des Zeugenbeweises offenen Voraussetzungen nicht frei vor.
\begin{tabproof}
 \proofstep{1}{M,s\models\exists x\,\varphi}{\rA}
 \proofstep{2}{\forall a\,(a\in D\to
      ((M,s[x\mapsto a]\models\varphi)\to Q))}{\rA}
 \proofstep{1}{\exists a\,
   (a\in D\land(M,s[x\mapsto a]\models\varphi))}
 {\FormulaRefAuto{B48Satisfaction}{1}}
 \proofstep{4}{b\in D\land(M,s[x\mapsto b]\models\varphi)}{\rA}
 \proofstep{4}{b\in D}{\rAEa{4}}
 \proofstep{4}{M,s[x\mapsto b]\models\varphi}{\rAEb{4}}
 \proofstep{2}{b\in D\to
      ((M,s[x\mapsto b]\models\varphi)\to Q)}{\rUE{2}}
 \proofstep{2,4}{(M,s[x\mapsto b]\models\varphi)\to Q}
 {\rRE{7,5}}
 \proofstep{2,4}{Q}{\rRE{8,6}}
 \proofstep{1,2}{Q}{\rEE{3,4:9}}
\end{tabproof}

\noindent
Die letzten beiden Saetze benutzen ausdruecklich die veraenderte
Belegung \(s[x\mapsto a]\). Daraus folgt noch nicht ohne Weiteres
die Korrektheit der syntaktischen Einsetzung \(\varphi[t/x]\).
Die folgenden Abschnitte beweisen dafuer das Koinzidenz- und
Substitutionslemma. Erst danach werden diese lokalen Schlussformen
auf syntaktische Herleitungen mit offenen Annahmen uebertragen.
